\section{Introducción}

La arquitectura Model-View-Controller (MVC) ha sido fundamental en el desarrollo de aplicaciones web y móviles desde su concepción en los años 70. Esta arquitectura separa las preocupaciones en tres componentes principales: el Modelo (Model), que maneja la lógica de negocio y los datos; la Vista (View), responsable de la presentación de la interfaz de usuario; y el Controlador (Controller), que gestiona la interacción entre el usuario y el sistema.

En el contexto actual, donde las aplicaciones web deben ser escalables, mantenibles y capaces de adaptarse a cambios rápidos, MVC ofrece una estructura clara que facilita el desarrollo colaborativo y la reutilización de código. Sin embargo, la implementación efectiva de MVC requiere no solo conocimientos técnicos sino también prácticas de ingeniería de software sólidas.

Este artículo presenta TuEvento, una plataforma completa de gestión de eventos desarrollada siguiendo los principios de MVC. El sistema permite a los usuarios descubrir, crear y gestionar eventos, incluyendo funcionalidades como autenticación, perfiles de usuario, compra de boletos y calificaciones. La implementación utiliza tecnologías modernas: Spring Boot para el backend, React para el frontend web y Expo/React Native para la aplicación móvil.

El desarrollo de TuEvento se basa en un proceso de ingeniería de software aplicado, utilizando modelado UML y generación automática de código. Esto permite evaluar empíricamente la calidad del código generado y validar la efectividad de la arquitectura MVC en contextos reales.

Nuestras contribuciones principales son: (1) una implementación completa de MVC en un sistema de gestión de eventos, (2) evaluación de calidad del código mediante métricas específicas para MVC, y (3) análisis de la aplicabilidad de MVC en el mercado actual de aplicaciones web.

El resto del artículo está estructurado como sigue: la Sección 2 revisa trabajos relacionados sobre MVC; la Sección 3 describe la metodología utilizada; la Sección 4 detalla la implementación; la Sección 5 presenta los resultados de la evaluación de calidad; la Sección 6 discute las implicaciones; y la Sección 7 concluye con lecciones aprendidas.